\chapter*{Annexes}
\addcontentsline{toc}{chapter}{Annexes}

\section*{Annexe A – Sujet du client SupHub reçu et traité dans ce caher des charges}

Nom du client : \textbf{SUPHUB} \\

\textbf{Contexte} : SUPHUB Solutions est une petite entreprise de conseil en technologies qui travaille fréquemment avec des équipes distribuées et des contributeurs externes sur divers projets open source. 

Actuellement, ils utilisent une combinaison d'outils disparates (tableurs pour le suivi de tâches, e-mails pour la communication, stockage cloud générique pour les fichiers) ce qui entraîne une perte d'efficacité, des doublons d'informations et une difficulté à suivre l'avancement global des projets. Ils souhaitent centraliser et optimiser leur gestion de projets. \\

\textbf{Besoin exprimé} : 

Nous avons besoin d'une plateforme web collaborative pour gérer nos projets open source de manière plus efficace. Cette plateforme doit nous permettre de :

Créer et gérer des projets (titre, description, statut, date de début/fin).
Assigner des membres aux projets avec des rôles (administrateur, développeur, relecteur).
Créer, attribuer et suivre des tâches (titre, description, responsable, date d'échéance, statut : à faire/en cours/terminé).
Permettre aux membres de commenter les tâches et les projets.
Partager et stocker des fichiers liés aux projets.
Visualiser l'avancement global d'un projet via un tableau de bord simple (ex: % de tâches terminées).
La solution doit être hébergée sur une infrastructure cloud pour assurer la disponibilité et la scalabilité, et doit être accessible depuis n'importe quel navigateur web.
Contraintes :

Une interface utilisateur intuitive et moderne.
Capacité à intégrer potentiellement avec des systèmes de contrôle de version (ex: GitHub/GitLab) dans le futur (non demandé pour cette version initiale).
Sécurité des données et gestion des accès robuste.